\documentclass[11pt]{article}
\usepackage[utf8]{inputenc}
\usepackage{amsmath,amssymb}
\usepackage[margin=1in]{geometry}
\usepackage{hyperref}
\emergencystretch=3em

\title{The mixed remainder with arbitrary shifts}
\author{Claude, with Daniel Murfet}
\date{2026-09-07}

\begin{document}
\maketitle
\sloppy

\begin{abstract}
The mixed term of the anchored decomposition (unit~88) was bounded by the constant-kernel block with
both exponents shifted by one. Here the bound is generalised to any amplitude dominated by
$Cu^{M_1}v^{M_2}(1+s)^De^{\beta sL}$ on the box: its block integral is at most
$Ce^{D^2/\beta}\cdot\mathrm{twoDGeneral}(\beta/2,2L,b,N^2,h_1+M_1,h_2+M_2,k_1,k_2)$, which is
$O\bigl(N^{-\min((h_1+M_1+1)/k_1,(h_2+M_2+1)/k_2)}(1+\log N)\bigr)$. With $M_1/k_1,M_2/k_2>\tau$
this is the remainder estimate $O(N^{-(p+\tau)}(1+\log N))$ of the rectangular face-jet decomposition
of the $d=2$ block. The polynomial envelope is absorbed by $(1+s)^D\le e^{D^2/\beta}e^{\beta s^2/2}$.
Zero \texttt{sorry}/\texttt{axiom}.
\end{abstract}

\section{The things formalised}
{\small
\begin{verbatim}
theorem one_add_pow_le_exp (β s : ℝ) (D : ℕ) (hβ : 0 < β) (hs : 0 ≤ s) :
    (1 + s) ^ D ≤ Real.exp ((D : ℝ) ^ 2 / β) * Real.exp (β / 2 * s ^ 2)
theorem twoDAmp_env_le (β b L C N : ℝ) (h₁ h₂ k₁ k₂ M₁ M₂ D : ℕ) (Φ : ℝ → ℝ → ℝ → ℝ)
    (hβ : 0 < β) (hN : 0 ≤ N) (hC : 0 ≤ C)
    (hΦ : Continuous fun x : ℝ × ℝ × ℝ => Φ x.1 x.2.1 x.2.2)
    (hbound : ∀ u ∈ Icc (0 : ℝ) b, ∀ v ∈ Icc (0 : ℝ) b, ∀ s, 0 ≤ s →
      |Φ u v s| ≤ C * (u ^ M₁ * v ^ M₂) * ((1 + s) ^ D * Real.exp (β * s * L))) :
    |twoDAmp β b N h₁ h₂ k₁ k₂ Φ|
      ≤ C * Real.exp ((D : ℝ) ^ 2 / β)
        * twoDGeneral (β / 2) (2 * L) b (N ^ 2) (h₁ + M₁) (h₂ + M₂) k₁ k₂
theorem twoDGeneral_shift_isBigO (β a b : ℝ) (h₁ h₂ k₁ k₂ : ℕ) (hβ : 0 < β) (hb : 0 < b)
    (hk₁ : 0 < k₁) (hk₂ : 0 < k₂) :
    (fun N : ℝ => twoDGeneral β a b (N ^ 2) h₁ h₂ k₁ k₂) =O[atTop]
      fun N : ℝ => N ^ (-min (((h₁ : ℝ) + 1) / k₁) (((h₂ : ℝ) + 1) / k₂)) * (1 + Real.log N)
theorem twoDAmp_env_isBigO ... :
    (fun N : ℝ => twoDAmp β b N h₁ h₂ k₁ k₂ Φ) =O[atTop]
      fun N : ℝ => N ^ (-min ((((h₁ + M₁ : ℕ) : ℝ) + 1) / k₁)
          ((((h₂ + M₂ : ℕ) : ℝ) + 1) / k₂)) * (1 + Real.log N)
\end{verbatim}
}

\section{Mathematics}

Pointwise on the box, $u^{h_1}v^{h_2}e^{-\beta s^2}|\Phi|\le Cu^{h_1+M_1}v^{h_2+M_2}e^{-\beta s^2}(1+s)^De^{\beta sL}$
with $s=Nu^{k_1}v^{k_2}$, and $(1+s)^D\le e^{Ds}\le e^{D^2/\beta}e^{\beta s^2/2}$ (the second step is
$Ds\le D^2/\beta+\beta s^2/2$, i.e. $(\beta s/2-D)^2\ge0$). The right side is then
$Ce^{D^2/\beta}$ times the integrand of the constant-kernel block $\mathrm{twoDGeneral}$ with
parameters $(\beta/2,2L,N^2)$ and exponents $(h_1+M_1,h_2+M_2)$, whose $N\to\infty$ behaviour is
known from the unconditional trichotomy of unit~47 (\texttt{twoDGeneral\_isBigO\_min}) after
composing with $N\mapsto N^2$ and converting $\log N^2=2\log N$.

\section{Paper-facing meaning}

In the rectangular decomposition $\Phi=\sum_{i<M_1}u^ia_i+\sum_{j<M_2}v^jb_j-\sum c_{ij}u^iv^j+R$
with $|R|\le Hu^{M_1}v^{M_2}$ (Astra~\#4(b)), the remainder contributes
$O(N^{-\min(p+M_1/k_1,p+M_2/k_2)}(1+\log N))$; choosing $M_1/k_1,M_2/k_2>\tau$ makes it
$O(N^{-(p+\tau)}(1+\log N))$, the remainder order of the $\tau$-cutoff theorem. Together with the
face-term expansions of unit~104 this is all the analysis needed; what remains is the decomposition
identity and the bookkeeping.

\section{Lean notes}

\begin{itemize}
\item \texttt{field\_simp} can close a goal outright; a trailing \texttt{ring} then errors with
"no goals".
\item After \texttt{Real.log\_pow} the factor \texttt{↑2} is a \texttt{Nat} cast, invisible to
\texttt{nlinarith}; \texttt{push\_cast} first.
\item \texttt{IsBigO.comp\_tendsto} with \texttt{tendsto\_pow\_atTop two\_ne\_zero} transports an
\texttt{atTop} estimate in $n$ to one in $N=\sqrt n$; \texttt{Real.rpow\_natCast} and
\texttt{Real.rpow\_mul} collapse $(N^2)^{-m/2}$ to $N^{-m}$.
\end{itemize}

\end{document}
